\section*{Заключение}
\addcontentsline{toc}{section}{Заключение}
В результате проведённого исследования подтверждена актуальность разработки
специализированного веб-приложения для автоматизации документооборота культурно-досуговых
учреждений. Анализ существующих систем электронного документооборота показал, что
коммерческие решения («1С:Документооборот», Directum, Alfresco) ориентированы
преимущественно на крупные организации и не отвечают специфике и бюджетным
ограничениям таких учреждений, как МУ ЦКиД «Факел» г.~Фрязино.

На основе полученных при разработке данных можно сформулировать следующие
\textbf{выводы}:
\begin{enumerate}
    \item Выбранный технологический стек (FastAPI, PostgreSQL, Jinja2)
    обеспечил оптимальный баланс между простотой разработки и производительностью системы:
    асинхронная обработка запросов FastAPI в сочетании с реляционной моделью данных
    PostgreSQL позволила реализовать надёжное хранение и маршрутизацию документов
    без избыточной технологической сложности.
    \item Реализация ролевой модели доступа (RBAC) ~\cite{rbac} с четырьмя ролями и механизмом
    JWT-аутентификации обеспечила разграничение прав между сотрудниками, руководителями
    структурных подразделений, художественным руководителем и ИТ-администратором,
    что соответствует реальной организационной структуре культурно-досугового учреждения.
    \item Внедрение конечного автомата для управления жизненным циклом мероприятий
    (На согласовании → Согласовано → Запланировано → Проведено) с ведением истории
    переходов повышает прозрачность процессов согласования и позволяет руководству
    оперативно контролировать ход выполнения культурных мероприятий.
    \item Механизм автоматической генерации отчётов в формате Excel (месячный,
    по сотрудникам, по структурным подразделениям) на основе библиотек pandas ~\cite{pandas} и openpyxl ~\cite{openpyxl}
    исключает ручной сбор статистики и существенно сокращает трудозатраты
    художественного руководителя при подготовке ежемесячной отчётности.
\end{enumerate}

\textbf{Репозиторий с кодом проекта:} https://github.com/xseniaKri/coursework\_rvp

\textbf{Хостинг проекта:} http://195.58.153.49:8000/

\textbf{Рекомендации и предложения.} Для практического внедрения разработанного
решения в деятельности МУ ЦКиД «Факел» рекомендуется развернуть приложение на
облачной платформе с настройкой резервного копирования базы данных PostgreSQL.
Это обеспечит сотрудникам бесперебойный доступ к системе при минимальных затратах
на ИТ-инфраструктуру и позволит отказаться от бумажного документооборота в рамках
основных процессов учреждения.

\textbf{Перспективы углубления темы.} В дальнейшем развитии системы целесообразно
рассмотреть реализацию модуля уведомлений (email или Telegram) для автоматического
информирования ответственных сотрудников об изменении статуса мероприятия,
а также расширение аналитического раздела за счёт интерактивных дашбордов
с визуализацией динамики мероприятий и посещаемости по периодам.

При разработке пояснительной записки к курсовому проекту строго соблюдались
требования~\cite{gost_report}, определяющие структуру и правила оформления
научно-исследовательских работ.